% !TEX root = machine_learning.tex
\chapter{Nearest-Neighbor Methods}
\label{chap:nearest.neighbors}

Unlike linear models,  nearest-neighbor methods  are completely non-parametric and
assume no regularity on the decision rule or the regression
function. In their simplest version, they require no training and rely
on the proximity of a new observation to those that belong to the
training set. We will discuss in this chapter how these methods are used for regression and classification, and study some of their theoretical properties.

\section{Nearest neighbors for regression}
\subsection{Consistency}
We let $\CR_X$ denote the input space, and $\CR_Y=\mR^q$ be the output space. We assume that a distance, denoted $\dist$ is defined on $\CR_X$. This means that $\dist: \CR_X\times \CR_X\to [0, +\infty]$ (we allow for infinite values) is a symmetric function such that $\dist(x, x') = 0$ if and only if $x=x'$ and, for all $x,x',x''\in \CR_X$
\[
\dist(x,x') \leq \dist(x,x'') + \dist(x'',x'),
\]
which is the triangle inequality.



Let $T = (x_1, y_1, \ldots, x_N, y_N)$ be the
training set. For $x\in\CR_X$, let 
\[
D_T(x) = (\dist(x, x_k), k=1, \ldots, N)
\]
 be the collection of all distances between $x$ and the inputs in the training set. 
We consider  regression estimators taking the form 
\begin{equation}
\label{eq:nn.general}
\hat f(x) = \sum_{k=1}^N W_k(x) y_k
\end{equation}
where $W_1(x), \ldots, W_N(x)$ is a family of coefficients, or weights, that only depends on $D_T(x)$.

We will, more precisely, use the following construction \cite{stone1977consistent}. Assume that a family of numbers $w_1 \geq w_2 \geq \cdots \geq w_N \geq 0$ is chosen, with $\sum_{j=1}^N w_j = 1$. Given $x\in \mR^d$ and $k\in\{1, \ldots, N\}$, we let 
$r_k^+(x)$ denote the number of indexes $k'$ such  that $\dist(x,x_{k'}) \leq \dist(x, x_k)$ and $r_k^-(x)$ the number of such indexes such  that $d(x,x_{k'}) < d(x, x_k)$. The coefficients defining  $\hat f$ in \cref{eq:nn.general} are then chosen as:
\begin{equation}
\label{eq:wj}
W_k(x) = \frac{\sum_{k' = r^-_k(x)+1}^{r^+_k(x)} w_{k'}}{r^+_k(x) - r^-_k(x)}.
\end{equation}
To emphasize the role of $(w_1, \ldots, w_N)$ is this definition, we will denote the resulting estimation as $\hf_w$.
If there is no tie in the sequence of distances between $x$ and elements of the training set, then $r^+_k(x) = r^-_k(x) + 1$ is the rank of $x_k$ when training data is ordered according to their proximity to $x$, and $W_k(x) = w_{r_k^+(x)}$. In this case, defining $l_1, \ldots, l_N$ such that $d(x, x_{l_1}) < \cdots < d(x, x_{l_N})$, we have
\[
\hf_w(x) = \sum_{j=1}^N w_j y_{l_j}.
\]
In the general case, the weights $w_j$ associated with tied observations are averaged.

If $p$ is an integer,  the $p$-nearest neighbor ($p$-NN) estimator (that we will denote $\hat f_p$) is associated to the weights $w_j = 1/p$ for $j = 1, \ldots, p$ and 0 otherwise. If there is no tie for the definition of the $p$th nearest neighbor of $x$, $W_k(x) = 1/p$ if $k$ is among the $p$ nearest-neighbors and $W_k(x) = 0$ otherwise, so that $\hf_p$ is the average of the output values over these $p$ nearest neighbors. If the $p$th nearest neighbors are tied, their output value is averaged before being used in the sum. For example, assume that $N = 5$ and $p=2$ and let the distances between $x$ and $x_k$ for $k=1, \ldots, 5$ be respectively $9, 3, 2, 4, 6$. Then $\hf_2(x) = (y_2+y_3)/2$. If the distances were $9, 3, 2, 3, 6$, then we would have $\hf_2(x) = (y_2+y_4)/4 + y_3/2$. 


When $\CR_X = \mR^d$ and  $d(x,x') = |x-x'|$, the following result is true.
\begin{theorem}[\cite{stone1977consistent}]
\label{th:stone}
Assume that
$\myE(Y^2) < \infty$. Assume that, for each $N$, a sequence $w^{(N)} = w^{(N)}_1 \geq  \cdots \geq w^{(N)}_N \geq 0$ is chosen with $\sum_{j=1}^N w^{(N)}_j =1$. Assume, in addition, that 
\begin{enumerate}[label=(\roman*)]
\item $\lim_{N\to\infty} w^{(N)}_1 = 0 $ 
\item $\lim_{N\to\infty} \sum_{j\geq \al N}w_j^{(N)} \to 0$, for some $\al\in (0,1)$. 
\end{enumerate}
Then the corresponding classifier $\hf_{w^{(N)}}$ converges in the $L^2$ norm to $\myE(Y\mid X)$: 
$$
\myE\left(|\hf_{w^{(N)}}(X) - \myE(Y\mid X)|^2\right) \to 0.
$$
For nearest-neighbor regression, (i) and (ii) mean that the number of nearest neighbors $p_N$ must be chosen such that $p_N \to \infty$ and $p_N/N \to 0$.
\end{theorem}

\begin{proof}
We give a proof under the assumption that $f: x \mapsto \myE(Y\mid X=x)$ is uniformly continuous and
bounded (one can, in fact, prove that it is always possible to reduce to this case).
%We are giving a partial proof of this result. Write
%$$
%F_N(X) = \sum_{i=1}^N W^{(N)}_{i}(X) Y_i
%$$
%with $W_{Ni}(x) = 1/|V_k(x)|$ if $i\in V_k(x)$ and 0 otherwise. 

To lighten the notation, we will not make explicit the dependency on $N$ in of quantities such as $W$ or $w$. One has
\begin{equation}
\label{eq:nn.proof}
\hf_w(X) - \myE(Y\mid X) =
\sum_{k=1}^N
W_{k}(X) (f(X_k) - f(X)) +
  \sum_{k=1}^N W_{k}(X) (Y_k - f(X_k)) 
\end{equation}
and the two sums can be addressed separately. 
%Starting with the second one, we have
%\[
%\abs{\sum_{i=1}^N
%W_{i}(X) (E(Y|X_i) - E(Y|X))} \leq \sum_{i=1}^N
%W_{i}(X) (E(Y|X_i) - E(Y|X))^2\,.
%\]

%
%The proof then consists in showing that
%$$
%E\left(\left|\sum_{i=1}^N W_{Ni}(X) (Y_i - E(Y|X_i))\right|^2\right)  \to 0
%$$ and
%$$
%E\left(\sum_{i=1}^N
%W_{Ni}(X) (E(Y|X_i) - E(Y|X))^2\right)  \to 0.
%$$
%Consider the second sum.
%The latter result proves the convergence of the second sum to 0 since
%We have
%$$
%\abs{\sum_{i=1}^N
%W_{i}(X) (E(Y|X_i) - E(Y|X))} \leq \sum_{i=1}^N
%W_{i}(X) (E(Y|X_i) - E(Y|X))^2
%$$
%from Cauchy-Schwartz inequality.
We start with the first sum  and write, by Schwartz's inequality:
$$
\left(\sum_k W_{k}(X)(f(X_k) - f(X))\right)^2 \leq \sum_k
W_{k}(X)(f(X_k) - f(X))^2.
$$
It therefore suffices to study the limit of 
$E(\sum_k W_{k}(X) (f(X_k)- f(X))^2$. Fix $\ep>0$. By assumption, there
exists $M,a>0$ such that $|f(x)| \leq M$ for all $x$ and $|x-y|\leq a
\Ria |f(x)-f(y)|^2 \leq \ep$.  Then
\begin{align*}
\myE\left(\sum_k W_{k}(X) (f(X_k)- f(X))^2\right) =& \myE\left(\sum_k W_{k}(X) (f(X_k)- f(X))^2\ \bfone_{|X_k-X|\leq
a}\right) \\
& + \myE\left(\sum_k W_{k}(X) (f(X_k)- f(X))^2\ \bfone_{|X_k-X|> a}\right) \\
&\leq \ep^2 + 4M^2 \myE\left(\sum_k W_{k}(X)  \bfone_{|X_k-X|> a}\right).
\end{align*}
Since $\ep$ can be made arbitrarily small, we need to show that, for any positive $a$, the second term in the upper-bound tends to 0  when $N\to\infty$.
We will use the following fact, which requires some minor measure theory argument to prove rigorously. Define
$$
S  = \defset{x: \forall \de > 0, \myP(|X-x|<\de) > 0}.
$$
This set  is called the support of $X$. Then, one can show that $\myP(X\in S) = 1$. This means
that, if $\tilde X$ is independent from $X$ with the same
distribution, then, for any $\de >0$,
$\myP(|X-\tilde X| <\de|X) > 0$ with probability one.
\footnote{This statement is proved  as follows (with the assumption that $X$ is Borel measurable). Let $S^c$ denote the complement of $S$. 
Then $S^c$ is open. 
Indeed if $x\not\in S$, there exists $\de_x>0$ such that, letting $B(x, \de_x)$ denote the open ball with radius $\de_x$, $\myP(X\in {B}(x, \de_x)) = 0$. Then $\myP(X\in B(x', \de_x/3)) = 0$ as soon as $|x-x'|<\de_x/3$, so that $B(x, \de_x/3)\sub S^c$.

If $K \subset S^c$ is compact, then  $K \subset \bigcup_{x\in K}  B(x, \de_x)$ and one can find a finite subset $M\subset K$ such that $K \subset \bigcup_{x\in M}  B(x, \de_x)$, which proves that $\myP(X\in K) = 0$. Since $\myP(X\in S^c) = \max_K \myP(X\in K)$ where the maximum is over all compact subsets of $S^c$, we find $\myP(X\in S^c) = 0$ as required.
}

Let $N_a(x) = |\defset{k: |X_k - x|\leq a} |$. We have, for all $x\in S$ and $a>0$,
and using the law of large numbers,
$$
\frac{N_a(x)}{N} = \frac{1}{N} \sum_{k=1}^N
\bfone_{|X_k-x|\leq a} \to P(|X-x| \leq a) > 0.
$$

If $|X-X_k| > a$, then $r^-_k(X) > N_a(x)$ so that 
\[
\sum_k W_{k}(X)  \bfone_{|X_k-X|> a} \leq \sum_{j\geq N_a(X)} w_j,
\]
and we have, taking $0 < \al < P(|X-x| \leq a)$,  
\[
\myE\left(\sum_{j\geq N_a(X)} w_j\right)\leq
\sum_{j\geq \al N} w_j + \myP(N_a(X) < \al N)
\]
and both terms in the upper bound converge to 0. This shows that the first sum in \eqref{eq:nn.proof} tends to 0.
\bigskip
%
%But, $\De_{k_N, N} \geq a$ implies $N_a(X)/N  \geq k_N/N$ so that
%$$P(\De_{N,k_N} > a) \leq P\left(\frac{N_a(X)}{N} \leq \frac{k_N}{N}\right).$$
%This tends to 0 as soon as $k_N/N\to 0$, since $N_a/N$ converges to a
%positive number with probability 1. \\

We now consider the second sum in \eqref{eq:nn.proof}. Let $Z_k = Y_k - \myE(Y\mid X_k)$. We have
$\myE(Z_k\mid X_k) = 0$ and $\myE(Z_k^2) <\infty$. We can write
\begin{eqnarray*}
\myE\left(\left|\sum_{k=1}^N W_{k}(X) Z_k\right|^2\right)  &=&\myE\left(\myE\left(\left|\sum_{k=1}^N W_{k}(X)
Z_k\right|^2\, \Big|\,X,X_1, \ldots, X_N\right)\right)\\
&=& \myE\left(\sum_{k=1}^N W_{k}(X)^2 \myE(Z_k^2 \mid X_k)\right) \\
&&+ \sum_{k\neq l=1}^N \myE(W_{k}(X)
W_{l}(X) \myE(Z_kZ_l \mid X_i, X_j))
\end{eqnarray*}
The cross products in the last term vanish because $\myE(Z_k\mid X_k) = 0$ and the samples are
independent. So it only remains to consider
$$
\myE\left(\sum_{k=1}^N W_{k}(X)^2 \myE( Z_k^2 \mid X_k)\right)
$$
The random variable $\myE(Z_k\mid X_k) = \myE(Y_k^2\mid X_k) - \myE(Y_k\mid X_k)^2$ is a fixed non-negative
function of $X_k$, that we will denote $h(X_k)$. 
%Since 
%$W_{N,i}(X) $ is either  0 { or } $1/|V_k(X)| \leq 1/k$, we can write
We have
$$
\myE\left(\sum_{k=1}^N W_{k}(X)^2 h(X_k)\right)\leq w_1
\myE\left(\sum_{k=1}^N W_{k}(X) h(X_i)\right)
$$
with $w_1\to 0$ and the proof is concluded by showing that $\myE\left(\sum_{k=1}^N W_{k}(X) h(X_k)\right)$ is bounded. 

Recall that the weights $W_k$ are functions of $X$ and of the whole training set, and we will need to make this dependency explicit and
write $W_i(X, \mT_X)$ where $\mT_X = (X_1, \ldots, X_N)$. Similarly, the ranks in \eqref{eq:wj} will be written $r^+_j(X, \mT_X)$ and $r^-_j(X, \mT_X)$.

Because $X, X_1, \ldots, X_N$ are i.i.d., we can switch the role of $X$ and $X_k$ in the $k$th term of the sum, yielding
\[
\myE\left(\sum_{k=1}^N W_{k}(X, \mT_X) h(X_k)\right) = \myE\left(\left(\sum_{i=1}^N W_{k}(X_k, \mT_X^{(k)})\right) h(X)\right)
\]
with $\mT_X^{(k)} = (X_1, \ldots, X_{k-1}, X, X_{k+1}, \ldots, X_N)$. We now show that 
$\sum_{k=1}^N W_{k}(X_k, \mT_X^{(k)})$ is bounded independently of $X, X_1, \ldots, X_N$. 

For this purpose, we group $X_1, \ldots, X_N$ according to approximate alignment with $X$. For $u\in\mR^d$ with $|u|=1$ and for $\de \in (0, \pi/4)$, denote by $\Ga(u, \de)$ the cone formed by all vectors $v$ in $\mR^d$ such that $\scp{v}{u} > |v|\cos\de $ (i.e., the angle between $v$ and $u$ is less than $\de$).  Notice that if $v, v'\in \Ga(u, \de)$, then $\scp{v}{v'} \geq \cos(2\de) |v|\,|v'|$ and if $|v'| \leq |v|$, then
\begin{equation}
\label{eq:nn.ineq}
|v|^2 - |v-v'|^2   = |v'|(2|v|\cos(2\de) - |v'|) > 0
\end{equation}
because $\cos(2\de) > 1/2$.


Fixing $\de$, let $C_d(\de)$ be the minimal number of such cones needed to cover $\mR^d$. Choosing such a covering $\Ga(u_1, \de), \ldots, \Ga(u_M, \de)$ where $M = C_d(\de)$, we define the following subsets of $\{1, \ldots, M\}$:
\begin{align*}
I_0 &= \defset{k: X_k = X}\\
I_q &= \defset{k\not \in I_0: X_k-X\in \Ga(u_q, \de)},\quad q=1, \ldots, M
\end{align*}
(these sets may overlap). We have
\[
\sum_{k=1}^N W_{k}(X_k, \mT_X^{(k)}) \leq \sum_{q=0}^M \sum_{k\in I_q} W_{k}(X_k, \mT_X^{(k)})
\]
If $k\in I_0$, then $r^-_k(X_k, \mT_X^{(k)}) = 0$ and $r^+_k(X_k, \mT_X^{(k)}) = c$ with $c = |I_0|$. This implies that, for $k\in I_0$, we have
$W_{k}(X_k, \mT_X^{(k)}) = \sum_{j=1}^c w_j/c$ and 
\[
\sum_{k\in I_0} W_{k}(X_k, \mT_X^{(k)}) = \sum_{j=1}^c w_j.
\]

We now consider $I_q$ with $q\geq 1$. Write $I_q = \{i_1, \ldots, i_r\}$ ordered so that $|X_{i_j} - X|$ is non-decreasing. If $j'<j$, we have (using \eqref{eq:nn.ineq})
$|X_{i_j}-X_{i_{j'}}| < |X-X_{i_j}|$. This implies that $r^-_{i_j}(X_{i_j}, \mT_X^{(i_j)}) \geq j-1$ and $r^+_{i_j}(X_{i_j}, \mT_X^{(i_j)}) - r^-{i_j}(X_{i_j}, \mT_X^{(i_j)}) \geq c+1$. Therefore,
\[
 W_{i_j}(X_{i_j}, \mT_X^{(i_j)}) \leq \frac1{c+1}{\sum_{j'=j}^{c+j} w_{j'}}
 \]
 and
\[
\sum_{k\in I_q} W_{k}(X_k, \mT_X^{(k)}) \leq \frac1{c+1}\sum_{j=1}^N \sum_{j'=j}^{c+j} w_{j'} = \frac1{c+1} \left(\sum_{j'=1}^c j'w_{j'} + (c+1) \sum_{j'=c+1}^N w_{j'}\right).
\]
This yields
\[
\sum_{k=1}^N W_{k}(X_k, \mT_X^{(k)}) \leq \sum_{j=1}^c w_j + C_d(\de) \left(\frac1{c+1}\sum_{j'=1}^c j'w_{j'} + \sum_{j'=c+1}^N w_{j'}\right)
\leq C_d(\de) +1.
\]
We therefore have 
\[
\myE\left(\sum_{k=1}^N W_{k}(X)^2 \myE( Z_k^2 \mid X_k)\right)\leq w_1(C_d(\de) +1) \myE(h(X)) \to 0,
\]
which concludes the proof.
 \end{proof}


Theorem \ref{th:stone} is proved in \citet{stone1977consistent} with weaker hypotheses allowing for more flexibility in the computation of distances, in which, for example, differences $X-X_i$ can be normalized by dividing them by a factor $\sig_i$ that may depend on the training set. These relaxed assumptions slightly complicate the proof, and we refer the reader to \citet{stone1977consistent} for a complete exposition. 


\subsection{Optimality}
The NN method can be shown to be optimal over some classes of
functions. Optimality is in the min-max sense, and works as follows. We assume that the regression function $f(x) = \myE(Y\mid X=x)$
belongs to some set $\CF$ of real-valued functions on $\mR^d$. Most of
the time, the estimation methods must be adapted to a given choice of
$\CF$, and various choices have arisen  in the literature: classes of
functions with $r$ bounded derivatives, Sobolev or related spaces,
functions whose Fourier transforms has given properties, etc.

Consider now an estimator of $f$, denoted $\hat f_N$, based on a
training set of size $N$. We can measure the error by, say:
$$
\norm{\hat f_N - f}_2 = \left(\int_{\mR^d} (\hat f_N(x) - f(x))^2
dx\right)^{1/2}
$$

Since $\hat f_N$ is computed from a random sample, this error is a
random variable. One can study, when $b_N\to 0$, the probability
$$
\myP_f\left(\|\hat f_N - f\|^2_2 \geq c b_N\right)
$$
for some constant $c$ and, for example, for the model: $Y = f(X) +
\text{noise}$. Here, the notation $\myP_f$ refers to the model assumption indicating the unobserved function $f$.

The min-max method considers the worst case and computes
$$
M_N(c) = \sup_{f\in \CF}  \myP_f\left(\|\hat f_N - f\|^2_2 \geq c b_N\right).
$$
This quantity now only depends on the estimation algorithm. One defines the notion of 
``lower convergence rate'' as a  sequence $b_N$
such that, for any choice of the estimation algorithm,  $M_N(c)$ can be found arbitrarily close to 1 (i.e., $\|\hat
f_N - f\|^2_2 \geq c b_N$ with arbitrarily high probability for all
$f\in \CF$), for
arbitrarily large $N$ (and for some choice of $c$). The mathematical
statement is
$$
\exists c > 0: \liminf_{N\to\infty} M_N(c) = 1.
$$
So, if $b_N$ is a lower convergence rate, then, for every estimator,
there exists a constant
$c$ such that the accuracy $cb_N$ cannot be achieved. 

On the other hand, one says that $b_N$ is an achievable rate of
convergence if there exists an estimator such that, for some $c'$,
$$
\limsup_{N\to\infty} M_N(c') = 0.
$$
This says that for large $N$, and for some $c'$, the accuracy is
higher than $c'b_N$ for the given estimator. Notice the difference: a
lower rate holds for all estimators, and an achievable rate for at
least one estimator.

The final definition of a min-max optimal rate is that it is both a
lower rate and an achievable rate (obviously for different constants
$c$ and $c'$). And an estimator is optimal in the min-max sense if it
achieves an optimal rate.


One can show that the $p$-NN estimator is optimal (under some assumptions
on the ratio $p_N/N$) when $\CF$ is the class
of Lipschitz functions on $\mR^d$, i.e., the class of functions such that
there exists a constant $K$ with
$$
|f(x) - f(y)|\leq K |x-y|
$$
for all $x,y\in\mR^d$. In this case, the optimal rate is $b_N =
N^{-1/(2+d)}$ (notice again the ``curse of dimensionality'': to achieve
a given accuracy in the worst case, the number of data points must grow exponentially with
the  dimension).

If the function class consists of smoother functions (for example,
several derivatives), the $p$-NN method is not optimal. This is
because the local averaging method is too crude when one knows already
that the function is smooth. But it can be modified (for example by
fitting, using least squares, a polynomial of some degree instead of computing an average)
in order to obtain an optimal rate.

\section{$p$-NN classification}

Let $(x_1,
y_1, \ldots, x_N, y_N)$ be the training set, with $x_i\in\mR^d$ and
$y_i\in \CR_Y$ where $\CR_Y$ is a finite set of classes. Using the same notation as in the previous section, define
\[
\widehat \pi_w(y|x) = \sum_{k=1}^N W_k(x) \bfone_{y_k=y}.
\]
Let the corresponding classifier be 
\[
\hf_w(x) = \argmax_{y\in \CR_Y} \widehat\pi_w(y|x).
\]

\Cref{th:stone} may be applied, for $y\in \CR_Y$, to
the function $f_y(x) = \pi(y\mid x) = \myE(\bfone_{Y=y} \mid X=x)$, which allows
one to interpret the estimator $\what\pi(y \mid x)$ as a nearest-neighbor
predictor of the random variable  $\bfone_{Y=y}$ as a function of $X$. We therefore obtain
the consistency of the estimated posteriors when $N\to\infty$ under the same assumption as those of \cref{th:stone}.
 This implies that, for large $N$, the classification will
be close to Bayes's rule.

An asymptotic comparison with Bayes's rule can already be made  with
$p=1$. Let $\hat y_N(x)$ be the 1-NN  estimator of $Y$
given $x$ and a training set of size $N$, and let $\hat y(x)$ be the
Bayes estimator. We can compute the Bayes error by
\begin{eqnarray*}
\myP(\hat y(X) \neq Y) &=& 1 - \myP(\hat y(X) = Y)\\
&=& 1 - \myE(\myP(\hat y(X) = Y|X))\\
&=& 1 - \myE(\max_{y\in\CR_Y} \pi(y|X))
\end{eqnarray*}

For the 1-NN rule, we have
\begin{eqnarray*}
\myP(\hat y_N(X) \neq Y) &=& 1 - \myP(\hat y_N(X) = Y)\\
&=& 1 - \myE(\myP(\hat y_N(X) = Y|X))\\
\end{eqnarray*}
Let us make the assumption that nearest neighbors are not tied (with probability one). Let $k^*(x, T)$ denote the index of  the nearest neighbor to $x$ in the training set $T$. We have
%Since $\hat y_N$ depends on the training set, the previous expectation
%can be written
\begin{eqnarray*}
\myP(\hat y_N(X) = Y\mid X) & = & \myE(\myP(\hat y_N(X) = Y\mid X, \mT))\\
& = & \myE\Big(\sum_{k=1}^N \myP( Y = Y_k\mid X, \mT) \bfone_{k^*(X, \mT) = k}\Big)\\
& = & \myE\Big(\sum_{k=1}^N \myP( Y = Y_k \mid X, X_k) \bfone_{k^*(X, \mT) = k}\Big)\\
& = & \myE\Big(\sum_{k=1}^N \sum_{g\in \CR_Y} \myP( Y = g,  Y_k=g\mid X, X_k) \bfone_{k^*(X, \mT) = k}\Big)\\
& = & \myE\Big(\sum_{k=1}^N \sum_{g\in \CR_Y} \pi( g\mid X) \pi(g\mid X_k) \chi_{k^*(X, \mT) = k}\Big)\\
& = & \myE\Big(\sum_{g\in \CR_Y} \pi( g\mid X) \pi(g\mid  X_{k^*(X,\mT)}) \Big)
\end{eqnarray*}

Now, assume the continuity of $x\mapsto \pi(g\mid x)$ (although the result
can be proved without this
simplifying assumption). We know that $X_{k^*(X,\mT)} \to X$ when
$N\to\infty$ (see the proof of \cref{th:stone}), which implies that
$\pi(g\mid X_{k^*(X,\mT)}) \to \pi(x\mid X)$ and at the limit
$$\myP(\hat y_N(X) = Y\mid X) \to \sum_{g\in \CR_Y} \pi( g\mid X)^2 .$$

This implies that the asymptotic 1-NN misclassification error is always smaller
than 2 times the Bayes error, that is
$$
1- \myE\Big(\sum_{g\in \CR_Y} \pi( g\mid X)^2 \Big) \leq 2(1- \myE(\max_g \pi(
g\mid X) ))
$$
Indeed, the left-hand term is smaller than $1- \myE(\max_g \pi(g|x)^2)$
and the result comes from the fact that for any $t\in\mR$. $1-t^2
\leq 2 - 2t$.

\begin{remark}
\label{rem:nn.comp}
Nearest neighbor methods may require large computation time, since, for a given $x$, the number of
comparisons which are needed is the size of the training set. However,
efficient (tree-based) search algorithms can be used in many cases to
reduce it to a logarithm in the size of the database, which is
acceptable. A reduction of the size of the training set by clustering also
is a possibility for improving the efficiency.

The computation time is also generally proportional to the dimension
$d$ of the input $x$. When $d$ is large, a reduction of dimension is
often a good idea. Principal components (see \cref{chap:dim.red}), or LDA directions (see \cref{chap:lin.class})
can be
used for this purpose.
\end{remark}

\section{Designing the distance}

\paragraph{LDA-based distance}
The most important factor in the design of a NN procedure probably is
the choice of the distance, something we have not discussed
so far. Intuitively, the distance should increase fast in the directions
``perpendicular'' to the regions of constancy of the class variables,
and slowly (ideally not at all) within these regions. The following construction uses  discriminant
analysis \cite{htf03}.

For $g\in \CR_Y$, let $\Sig_g$ be the covariance matrix in class $g$, and $\Sig_w =
\sum_{g\in \CR_Y}\pi_g \Sig_g$ be the within-class variance, where $\pi_g$ is the
frequency of class $g$. Let $\Sig_b$ denote the between-class covariance matrix (see \cref{sec:lda}).

For $x\in\mR^d$, define the spherized vector
$x^* = \Sig_w^{-1/2} x$. The between-class variance computed for spherized
data is $\Sig_b^* = \Sig_w^{-1/2} \Sig_b \Sig_w^{-1/2}$. A direction is discriminant if it
is close to the principal eigenvectors of $\Sig_b^*$. This suggests the
introduction of the norm
$$
|x|_*^2  = (x^*)^T \Sig_b^* x^* = x^T \Sig_w^{-1/2}
(\Sig_w^{-1/2}\Sig_b\Sig_w^{-1/2})\Sig_w^{-1/2}x = x^T \Sig_w^{-1}\Sig_b\Sig_w^{-1}x.
$$
This replaces the standard Euclidean norm (the method can be made more
robust by adding $\ep\Id[d]$ to $\Sig_b^*$.)

\paragraph{Tangent distance}
Designing the distance, however, can sometimes be based on {\it a
priori} knowledge on some invariance properties associated with the classes. A successful example comes from
character recognition, where it is known that transforming images by
slightly rotating, scaling, or translating the character should not
change its class. This corresponds to the following general framework.

For each input $x\in\mR^d$, assume that one can make small transformations without
changing the class of $x$. We model these transformations as
parametrized functions $x \mapsto x_\th = \phi(x, \th)\in\mR^d$, such that $\phi(x,
0) = x$ and $\phi$ is smooth in $\th$, which is a $q$-dimensional parameter. The assumption is that $\phi(x, \th)$
and $x$ should be from the same class, at least for small $\th$. This
will be used to improve on the Euclidean distance on $\mR^d$.

Take $x, x'\in\mR^d$. Ideally, one would like to use the distance $D(x,x') =
\inf_{\th, \th'} \dist(x_\th, x_{\th'})$ where $\th$ and $\th'$ are
restricted to a small neighborhood of 0. A more tractable
expression can be based on first-order approximations
\begin{align*}
&x_\th \simeq x + \nabla_{\th} \phi(x,0) u = x + \sum_{i=1}^q u_i \prt_{\th_i}\phi(x,0)\\
\text{and}\quad & x'_\th \simeq x' + \nabla_{\th} \phi(x',0) u' =x' + \sum_{i=1}^q u'_i \prt_{\th_i}\phi(x',0)
\end{align*}
yielding the approximation (also called the tangent distance)
$$
D(x, x')^2 \simeq \inf_{u,u'\in\mR^q} \norm{x - x' + \nabla_{\th} \phi(x,0) u - \nabla_{\th} \phi(x',0)
u'}^2.
$$
The computation now is a simple least-squares problem, for which the solution is
given by the system
$$
\begin{pmatrix} \nabla_\th \phi(x,0)^T\nabla_\th \phi(x,0) & -\nabla_\th \phi(x,0)^T\nabla_{\th}\phi(x',0)\\
-\nabla_{\th}\phi(x',0)^T\nabla_\th \phi(x,0) & \nabla_{\th}\phi(x',0)^T\nabla_{\th}\phi(x',0)\end{pmatrix} 
\begin{pmatrix}u\\ v\end{pmatrix} = 
\begin{pmatrix} \nabla_\th \phi(x,0)^T (x'-x) \\ \nabla_{\th} \phi(x',0)^T (x-x')\end{pmatrix}.
$$
A slight modification, to ensure that the norms of $u$ and $u'$ are
not too large, is to add a penalty $\la (|u|^2 + |u'|^2)$, which
results in adding $\la \Id[q]$ to the diagonal blocs of the above matrix.

